\section{Finite optimization and exact sampling}
\label{sec:numerics}

The optimization step is a finite implementation of the conditional-gradient
method of Frank and Wolfe~\cite{FrankWolfe1956}. We use the fixed-step
form analyzed by Jaggi~\cite[Algorithm~2 and Theorem~1]{Jaggi2013}, applied
to a smoothed entropy objective. Our linearization has a fixed error floor,
whose effect is included in the recurrence below. The rational arithmetic
and exact sampling are also part of the running-time analysis.

The online rule in Section~\ref{sec:bands} asks for a mixture of band
actions whose mean vector is close to an entropy maximizer.  Its finite
implementation separates recursive evaluation from numerical work.
First gather the earlier-index child values, once for each distinct
coordinate and history index occurring in the band.  The procedure below
then reuses this table while scoring the listed actions and constructing
a rational mixture.  It returns one exactly sampled action; only the
evaluation of that action requests current-index child values.  Thus the
mean vector is used to analyze the draw, whereas the returned refinement
value uses the cost of the single sampled action.

Fix an unskipped band $\mathcal C=\mathcal C_{b,r}$, and write
$d=d_{\mathcal C}$ for its number of coordinates.  Its vertices
$\rho(\alpha)$ are probability vectors, with child-block masses
$h_i/h_\Sigma$.  Let $v\in\mathbb Q^d$ be the coefficient vector
\[
 v=\eta\left(\sum_{j<t}u_j+\widehat g_t\right).
\]
Thus the desired vector $\rho_t$ maximizes
\[
 \mathcal F(\rho)=\mathcal H(\rho)+\langle v,\rho\rangle
 \quad\text{over}\quad
 \mathcal Q=\operatorname{conv}\{\rho(\alpha):\alpha\in\mathcal C\}.
\]
Every coordinate occurs in a listed vertex.  Also $h_i\ge b/M>0$, so all
the prescribed block masses are positive, even on an execution with an
inaccurate seed.

\begin{lemma}[A finite entropy optimizer]
\label{lem:finite-optimizer}
Let $\mathcal Q$ be the convex hull of a nonempty finite list of rational
probability vectors in $\mathbb R^d$, and suppose that every coordinate
is positive in at least one listed vector.  Given a rational vector $v$
and an integer $H\ge4$, put $G_0=H+d+2$.  A rational algorithm constructs
a mixture of the listed vectors whose mean $\rho$ satisfies
\[
 \norm{\rho-\rho_t}_1\le H^{-3},
\]
where $\rho_t$ is the unique maximizer of
$\mathcal H(\rho)+\langle v,\rho\rangle$ on $\mathcal Q$.
The number of bit operations and the stored mixture size are polynomial
in $G_0$, the list size, and the maximum numerator or denominator bit
length of an input rational.
The mixture can be sampled exactly using fair random bits, with expected
cost polynomial in the same quantities and with a fixed space bound.
\end{lemma}

\begin{proof}
There is a feasible vector positive in every coordinate, obtained by
averaging the list.  The maximizer is positive as well: at a zero
coordinate, moving a small distance towards that average produces a
positive entropy contribution of order $s\ln(1/s)$, which dominates
the terms of order $s$ from the other coordinates and the linear
objective.  Strict concavity gives uniqueness.

Set $\sigma=G_0^{-50}$ and smooth the objective by replacing
$-x\ln x$ with $-(x+\sigma)\ln(x+\sigma)$ in each coordinate.  Denote
the result by $\mathcal F_\sigma$.  Integrating the derivative of
$x\ln x$, including at zero, gives
\begin{equation}
 \sup_{\rho\in\mathcal Q}
 \abs{\mathcal F_\sigma(\rho)-\mathcal F(\rho)}
 \le 4d\sqrt\sigma.
 \label{eq:entropy-smoothing}
\end{equation}
The gradient of $\mathcal F_\sigma$ has Euclidean Lipschitz constant
at most $1/\sigma$.  The diameter of $\mathcal Q$ is at most $2$.

Start at a listed vertex.  At step $j\ge0$, approximate every coordinate
of the gradient to absolute error at most $e_0=G_0^{-20}$.  Enumerate
the list and choose a vertex $v_j$ maximizing its scalar product with
this approximate gradient, resolving ties by a fixed order.  Update
\[
 \rho^{(j+1)}=(1-\xi_j)\rho^{(j)}+\xi_j v_j,
 \qquad \xi_j=\frac{2}{j+2}.
\]
Here $v_j$ denotes a selected vertex, whereas $v$ without a subscript
is the fixed linear coefficient vector.  The difference between two
probability vectors has $\ell_1$ norm at most $2$.  Consequently the
chosen vertex loses at most $2e_0$ in the true linearization, even
relative to a feasible mixture.  If $D_j$ is the gap from the maximum
of $\mathcal F_\sigma$, concavity and the gradient bound give
\begin{equation}
 D_{j+1}\le(1-\xi_j)D_j+2e_0\xi_j+
                  \frac{2\xi_j^2}{\sigma}.
 \label{eq:finite-optimizer-recursion}
\end{equation}
Since $\xi_0=1$, this estimate first gives
$D_1\le2e_0+2/\sigma$, independently of the size of the linear
coefficient vector.  Induction then gives
\[
 D_j\le2e_0+\frac{8}{\sigma(j+1)}\qquad(j\ge1).
\]
Run $K_{\mathrm{opt}}=G_0^{80}$ steps.  By
\eqref{eq:entropy-smoothing}, the resulting true objective gap is at
most
\[
 2G_0^{-20}+8G_0^{-30}+8dG_0^{-25}
 \le12G_0^{-20}.
\]

The Hessian of negative entropy is
$\operatorname{diag}(1/\rho_k)$, which dominates the identity on the
probability simplex.  The first-order inequality at the positive
maximizer, followed by this Hessian bound, therefore shows that its
objective gap at any feasible $\rho$ is at least
$\frac12\norm{\rho-\rho_t}_2^2$.  The statement at coordinates equal
to zero follows by continuity.  Hence
\[
 \norm{\rho-\rho_t}_1
 \le\sqrt d\,\norm{\rho-\rho_t}_2
 \le\sqrt{24d}\,G_0^{-10}\le H^{-3}.
\]

We next give an exact representation of every iterate.  After $K\ge1$
steps, the weight of the vertex selected at step $j$, where $0\le j<K$,
is
\begin{equation}
 \frac{2}{j+2}\prod_{k=j+1}^{K-1}\frac{k}{k+2}
 =\frac{2(j+1)}{K(K+1)}.
 \label{eq:optimizer-mixture-weights}
\end{equation}
The initial vertex has zero remaining weight because $\xi_0=1$.
Store the sequence of selected vertices, allowing repetitions.  The
displayed weights sum to one and give the iterate exactly.  A common
denominator for the input vertices, multiplied by $K(K+1)$, is also a
common denominator for its coordinates.  One may obtain the input
denominator by multiplying all input denominators, whose total bit
length is at most the input size.  Therefore the iterates can be
maintained by rational arithmetic of polynomial bit cost, without
accumulating the denominators of successive numerical approximations.

It remains to calculate the approximate gradients using finite
rational operations.
For $r=x+\sigma\in[\sigma,1+\sigma]$, put
$b_0=(r-1)/(r+1)$.  Then $\abs{b_0}\le1-\sigma$, and
\[
 \ln r=2\sum_{k\ge0}\frac{b_0^{2k+1}}{2k+1}.
\]
Truncation after $K_{\log}=G_0^{60}$ terms has absolute error at most
\[
 \frac{2\abs{b_0}^{2K_{\log}+1}}
      {(2K_{\log}+1)(1-b_0^2)}
 \le
 \frac{2e^{-2K_{\log}\sigma}}
      {(2K_{\log}+1)\sigma}
 <G_0^{-20}.
\]
All terms are rational.  The linear part of the gradient and its
constant entropy derivative term are computed exactly.  Thus these
finite sums provide the required gradient precision; choosing a
vertex requires only rational comparisons.

Computing each rational logarithm sum by direct integer arithmetic
also has polynomial bit cost: the lengths
of the powers of a rational and of the products of the odd integer
denominators are polynomial in $K_{\log}$ and the input bit length.
These logarithmic approximations only select vertices; they are not
returned as recursive table values.

The same representation gives an exact sampler.  Put
$K=K_{\mathrm{opt}}$ and sample a uniform integer
$R\in\{1,\ldots,K(K+1)/2\}$ by drawing uniformly from a range whose
size is the least power of two at least $K(K+1)/2$, and rejecting
out-of-range values.  Each trial
succeeds with probability greater than $1/2$.  Choose the smallest
$j$ with $(j+1)(j+2)/2\ge R$, by integer binary search, and return
$v_j$.  Equation~\eqref{eq:optimizer-mixture-weights} proves exactness,
including when the same vertex was selected more than once.  The
integer has $O(\log K)$ bits.  Rejection reuses a fixed buffer, and
its expected bit cost is polynomial.  This completes the construction.
\end{proof}

\begin{corollary}[Local cost of the online rule]
\label{cor:online-local-cost}
For a band in the edit-distance algorithm, the numerical mixture in
Section~\ref{sec:bands} can be constructed and sampled with local
expected bit cost $2^{o(H)}$, using only already gathered child values
of earlier indices.  The storage used by this calculation is
$2^{o(H)}$ words.
\end{corollary}

\begin{proof}
At a wide-action endpoint there are $O(P+1)$ possible grid positions.
Thus the full wide list has at most $(C(P+1))^{2M}$ members, for an
absolute constant $C$, and the band is a sublist.  Also
$d\le M(21P)^2$, so $G_0=H^{O(p_*+1)}$.
For these vertices, use the common denominator $Mh_\Sigma$:
every $Mh_i$ is an integer because the dyadic scales satisfy
$h_i\ge b/M$ and $b\ge1$.  The guard $h_\Sigma\le64Fb$
bounds this denominator's bit length by $O(H)$, without multiplying
denominators across the action list.
Lemma~\ref{lem:finite-optimizer} gives a polynomial number of bit
operations in these quantities and in the bit length of each input
rational.  The size bounds proved in Section~\ref{sec:work} make the last
quantity $H^{O(p_*+1)}$ too.  The total is therefore at most
\[
 2^{O(M(p_*+1)\ell)}H^{O(p_*+1)}=2^{o(H)},
\]
because $M\le H^{1/20}$, $\ell=\Theta(\log H)$, and
$p_*+1=O((\log\ell)^{1/4}+1)$.
The same bounds apply to storage, with a reused rejection buffer.
Every list score and every numerical iteration uses the stored child
inputs, rather than a new recursive calculation for each action.
\end{proof}
